\cotefile{93}
\scanpage{135}
la quelle extension s'identifie à $\mathrm{P}^{1}_{X/S}(\omega)$
par sa structure d'extension naturelle de
$\omega$ par $\underline{\omega}\otimes\deleted{\unsure{\underline{\omega}}}\Omega^1_{X/S}\simeq\underline{\omega}\otimes\underline{\omega}=\underline{\omega}^{2}$) s'identifie
à l'op. diff. de degré $1$ (rel. à $S$) universel
sur $\underline{\omega}$
\[
C\vert\underline{\omega}=d^{1}_{\underline{\omega},X/S}:
\underline{\omega}\longrightarrow{\unsure{\mathscr P|\mathrm P|\mathcal P}}^{1}_{X/S}(\underline{\omega}).
\]

L'indétermination, \inserted{donc,} en cette construction,
au lieu d'être dans $\Gamma(\underline{\omega}\otimes\mathscr J)$
\deleted{\unsure{est}} (opérant sur $\mathscr J$ par crochet, à valeurs
dans $\mathscr J\otimes\underline{\omega}$) est \deleted{dans} \inserted{une section de}
$\underline{\omega}\otimes\mathrm{Fil}_0\simeq\underline{\omega}\otimes\underline{\omega}\simeq
\underline{\omega}^{\otimes 2}$ -- on retombe sur ses pattes !

\smallskip\centerline{\rule{1.2cm}{0.35pt}}\smallskip

Ceci ``rappelé'', on voit donc que si $S=\operatorname{Spec}\mathbb C$,
\deleted{\unsure{la|donc}} $X$ courbe algébrique lisse $/\mathbb C$, alors
la \deleted{\unsure{argument}}\inserted{construction} \uline{transcendante} \deleted{\unsure{appliquée}}\inserted{indiquée}
plus haut \deleted{\unsure{nous fournit de|nous fournit des}}\inserted{s'interprète comme la} construction
d'une connexion projective (canonique)
sur $X/\mathbb C$, i.e. d'une section canonique
du torseur ${\unsure{\mathscr C|C|\mathcal C}}_{X/\mathbb C}$ \inserted{sous $\underline{\omega}_{X/\mathbb C}^{\otimes 2}=\underline{\omega}^{2}$} des connexions
projectives sur $X$. \inserted{On se demande quand} cette connexion
projective \deleted{\unsure{est-elle|est}} de plus ``régulière à
l'infini'' i.e. la connexion correspondante sur
$\mathbb P_X\simeq\mathbb P(\mathscr J/\mathscr J^{3})$ ($\mathscr J$ Idéal
d'augmentation de $\mathscr P^\infty_{X/\mathbb C}$) \deleted{étant}\inserted{est} régulière
à l'infini. Sauf erreur, cela doit signifier

% margin note written sideways
\margin{par un argument analogue on voit que sur une courbe elliptique privée d'un pts, la connex. projective canonique n'est pas régulière, \unsure{Voir|Vois} (\unsure{autrement} elle ne \unsure{dépendrait} pas du dit trou, par raison d'\unsure{homogénéité} -- \unsure{trop rapide|très rapide} !)}
% slanted
\margin{Ce n'est pas le cas sur ${\unsure{E|\mathbb E}}_{\mathbb C}^{1}$ à $\unsure{0|1}$ \unsure{trou}, \unsure{ni} sur $\mathbb P^1_{\mathbb C}$ privée de 3 pts, \unsure{enfin} (car sur $\mathbb P^1_{\mathbb C}$ \unsure{privé de} $\nu\leq3$ pts, il n'y a au \unsure{plus} qu'une seule conn. proj. régulière (à l'$\infty$ !))}
% archive sheet number: 134
\scanpage{136}
que cette section du torseur (qui est
en fait défini sur $\widehat X$, compactifiée
non singulière naturelle de $X$, \unsure{\}|)} soit
${\unsure{\mathscr C|C|\mathcal C}}_{\widehat X/\mathbb C}$) n'admet ``à l'infini'' $\widehat X-X$
que des pôles simples au plus.

Si $D$ est le diviseur correspondant
\inserted{sur $\widehat X$} à multiplicités toutes $1$, de support $\widehat X-X$,
on aura comme groupe d'indétermination
des conn. proj. rég. à l'infini sur $\widehat X$
\[
 H^0(\widehat X,\underline\omega_{\widehat X}^{\otimes2}(D))
\]
qu'on peut relier ainsi aux invariants classiques
de la courbe complète $\widehat X$ (en écrivant
\deleted{\unsure{maintenant}}\inserted{d'abord} $\deleted{\underline{\omega}^{2}}\ \underline\lambda\deleted{\unsure{\otimes2}}$ au lieu de $\underline{\omega}^{\otimes2}_{\widehat X/\mathbb C}$)
\[
\deleted{0\longrightarrow\underline{\omega}^{2}\longrightarrow
 \underset{\deleted{\unsure{\Theta}}}{\inserted{\underline\lambda(D)}}\longrightarrow
 \underline{\omega}^{2}(D)/\underline{\omega}^{2}\longrightarrow0}
\]
\[
0\longrightarrow\underline\lambda\longrightarrow\underline\lambda(D)
\longrightarrow\underline\lambda(D)/\underline\lambda\longrightarrow0
\]
déduite par $\otimes\,\underline\lambda$ de

\diagram{\begin{tikzcd}[ampersand replacement=\&,column sep=small,row sep=tiny]
0\arrow[r]\&\underline O_{\widehat X}\arrow[r]\&
\underline O_{\widehat X}(D)\arrow[r]\&
\underline O_{\widehat X}(D)/\underline O_{\widehat X}\arrow[r]\&0\\
\&\&\underbrace{\mathscr J_D^{-1}}_{\text{Idéal de }D}\arrow[u,equal]\&
(\mathscr J_D/\mathscr J_D^2)^{-1}\arrow[u,"\unsure{\wr|\simeq}" description]\\
\&\&\&\deleted{\unsure{(\underline\omega|D)^{-}}}\arrow[u,"\unsure{\wr|\simeq}" description]\\[-1ex]
\&\&\&(\underline\omega^{-1}|D)
\end{tikzcd}}

d'où
\[
\underline\lambda(D)/\underline\lambda\simeq(\underline\lambda\,\underline\omega^{-1})|D
\]
et la suite exacte

\diagram{\begin{tikzcd}[ampersand replacement=\&,column sep=tiny,row sep=small]
0\arrow[r]\&H^0(\widehat X,\underline\lambda)\arrow[r]\&
H^0(\widehat X,\underline\lambda(D))\arrow[r]\&
H^0(D,\underline\lambda\,\underline\omega^{-1}|D)\arrow[d]\\
\&\&\&H^1(\widehat X,\underline\lambda).
\end{tikzcd}}

(ici $\underline\lambda=\underline\omega_{\widehat X/\mathbb C}^{\otimes2}\simeq\underline\omega^{2}$

\inserted{On suppose $X$ connexe,} \textcircled{a} Si $\widehat X$ est de genre $g\geq2$, $H^1(\widehat X,\underline\lambda)=0$
(c'est le dual de $H^0(\widehat X,\underline\omega^{-1})=0$) donc on
trouve

\diagram{\begin{tikzcd}[ampersand replacement=\&,column sep=tiny,row sep=small]
0\arrow[r]\&H^0(\widehat X,\underline\omega^2)\arrow[r]\&
H^0(X,\underline\omega^2)^{\text{rég}}\arrow[r]\&
H^0(D,\underline\omega|_{\deleted{\unsure{X}}D})\arrow[r]\&0\\
\&\unsure{:=|=}\Omega^1_{M_g/\mathbb C}\,(s)\arrow[u,"\unsure{\wr|\simeq}" description]\&
\Omega^1_{M_{g,\nu}/\mathbb C}\,(s)\arrow[u,"\unsure{\wr|\simeq}" description,"{\text{can. ?}}"']\&
\text{rang }\nu=\operatorname{card}\deleted{\unsure{D/}}\inserted{\operatorname{supp}D}
\end{tikzcd}}

\margin{rang \deleted{\unsure{g}} $3g-3$}
$M_g$ variété modulaire des courbes de genre $g$
\inserted{lisses complètes} avec $M_{g,\nu}$ variété modulaire \unsure{pour}
pointées à $\nu$ trous
% archive sheet number: 135
\scanpage{137}
\textcircled{b} $g\leq1$, i.e. $g=0,1$. Est-ce que
\[
H^1(\widehat X,\underline\lambda)\longrightarrow H^1(\widehat X,\underline\lambda(D))
\quad\text{est inj. ?}
\]
ou encore

\diagram{\begin{tikzcd}[ampersand replacement=\&,column sep=large,row sep=small]
H^0(\widehat X,\underline\omega^{-1})\&H^0(\widehat X,\underline\omega^{-1}(-D))\arrow[l]\\
\underline{\operatorname{Lie}}\,\underline{\operatorname{Aut}}(\widehat X)\arrow[u,equal]\&
\underline{\operatorname{Lie}}\,\underline{\operatorname{Aut}}(X)\arrow[u,equal]\\
{\underline t_{\widehat X/\mathbb C}\quad\text{fibré tangent}}\arrow[uu,equal]\&
{\underline t_{\widehat X/\mathbb C}(-D),\quad\text{champs \unsure{tangts|tangents} s'annulant sur }D}\arrow[uu,equal]
\end{tikzcd}}

\deleted{injectif} surjectif ? Si son conoyau \unsure{$\rightarrow$}\inserted{\unsure{sans flèche|sans plus|sans pôles}}
$Q(X)$, on aura
\[
0\longrightarrow H^0(\widehat X,\underline\omega^2)\longrightarrow
H^0(X,\underline\omega^2)^{\mathrm{rég}}\longrightarrow
H^0(D,\underline\omega|D)\longrightarrow\check Q(X)\longrightarrow0
\]

\uline{Cas} $g=1$, on \deleted{\unsure{dim}} $\operatorname{rg} H^0(\widehat X,\underline t_{\widehat X/\mathbb C})=\deleted{\unsure{\operatorname{rg}H^1\mathbb C}}1$
et si $D\ne0$ i.e. $X\ne\widehat X$ i.e. $X$ non complète,
on a $Q(X)\simeq H^0(\widehat X,\underline\omega^{-1})\simeq t_{\widehat X}$
(\deleted{\unsure{espace tangent}}\inserted{Algèbre de Lie} de
la Jacobienne) d'où $\check Q(X)=\omega_{\widehat X}=\check t_{\widehat X}$ (formes
diff. inv. sur jacobienne). \unsure{Interprét.|Interprétation} la plus \unsure{nette} :
on regarde les Pb modulaires des \uline{courbes}
\uline{elliptiques} (à \uline{pt marqué}) avec $\nu$ pts
distincts \uline{de somme nulle} \deleted{\unsure{et celà est que si elles sont}}\inserted{(ou ce qui revient au même,}
courbes elliptiques \unsure{princ.} homogènes et
$\nu$ pts \deleted{marqués}\inserted{(non ordonnés)} distincts et un pt choisi
de \unsure{tel} que ce soit la somme desdits \ldots.

\uline{Cas} $g=0$ \inserted{Alors $H^0(\widehat X,\underline t_{\widehat X/\mathbb C})=0$,}
$H^0(\widehat X,\underline\omega^{-1}_{\widehat X/\mathbb C})$ est de dim $3$, et
suivant les cas

\diagram{$\begin{array}{|ll}
\nu=1&\dim H^0(\widehat X,\underline t_{\widehat X/\mathbb C}(-D))=2\\
\nu=2&\rule{3.7cm}{0.35pt}=1\\
\nu\geq3&\rule{3.7cm}{0.35pt}=0
\end{array}$}

donc $Q(X)$ est respectiv$^{\mathrm{t}}$ de dim $1,2$ ou $3$.
Donc pour $\uline{0\leq\nu\leq3}$, on aura $\dim\check Q(X)=\operatorname{card}\operatorname{supp}D
=\dim H^0(D,\underline\omega|D)$
% archive sheet number: 136
\scanpage{138}
donc $H^0(X,\underline\omega^2)^{\mathrm{rég}}=0$ :
\uline{Il n'y a qu'une seule connexion projective} \inserted{régulière à l'$\infty$} \uline{sur $X$} si
\uline{$X$ est de genre $0$ et à $\nu\leq3$ trous}. Et
par contre $\nu\geq4$, l'espace $H^0(X,\underline\omega^2)^{\mathrm{rég}}$ est
de dim $\nu-3$ -- c'est encore la dim de
l'espace modulaire pour la droite projective
à $\nu$ trous, et sans doute on trouve
encore un isom can $H^0(X,\underline\omega^2)^{\mathrm{rég}}
\simeq\Omega^1_{M_{0,\nu}}(s)$ \ldots\ Bien entendu, dans le
cas $\nu\leq3$, la seule et unique connexion
projective sur $X$ est celle induite par
l'unique conn. projective de $\widehat X$ (et leur
description en \unsure{termes|tant} de conn. n'a rien
de transcendant !). Dans le cas de $\nu\leq3$,
% slanted
\margin{\unsure{Les|ses} \unsure{résultats} restent bien \unsure{entendu} \unsure{vrais} sur une base $S$ \unsure{quelconque} \unsure{au lieu de} \unsure{$\mathbb C$}.}
\deleted{\unsure{les}} la connexion projective \inserted{unique} \deleted{est d'ailleurs}
$\mathbb P_X=\widehat X\times X$, muni de la section diagonale \inserted{\unsure{(évidente)|(d'ailleurs)}} et
de la connexion triviale. Dans le cas
de $X$ quelconque, \inserted{(de genre $0$)} cette
connexion projective est tjrs définie, et
fournit une \uline{section canonique}.
% archive sheet number: 137
\scanpage{139}
\uline{Prop} \inserted{On suppose $2$ inv. sur $S$.}
Soit $H$ schéma en groupes lisse sur $S$,
de dim rel. $1$, et $X$ un torseur sous $H$.
Soit $(\unsure{\mathscr J|\mathscr V|\mathcal V},\unsure{\mathscr L|L|\mathcal L},\deleted{\unsure{O|\underline\omega}})$ l'Alg. de Lie de rang loc. $3$ canonique
sur $X$, \inserted{$(\mathbb P,\sigma)$ le fibré proj. à section associé,}
$H$ y opère par translation et le descend
par
\[
\underbrace{\underline\omega_0}_{\omega_0}\subset
\underbrace{L_0}_{\underline O_S}\subset
\underbrace{\mathscr J_0}_{\omega_0^{-1}}\ \text{sur } X,\qquad
(\mathbb P_0,\sigma_0)\quad \mathbb P_0=\mathbb P(L_0).
\]
\[
L_0\deleted{\unsure{\simeq|=}}\simeq(\underline O_H/\underline{\mathscr J}^3)^+,\quad
\underline{\mathscr J}\ \text{id. d'augm de }H
\]

Alors les conn. affines \inserted{sur $X$} invariantes sous $H$
sont en corr. $1$-$1$ avec \inserted{les sections d'}un certain sous
torseur \deleted{\unsure{sur}} de $\mathscr J_0\otimes\underline\omega_0$ sous $\underline\omega_0\otimes\underline\omega_0=\underline\omega_0^2$,
se projettant sur la section $1$ de $\underline\omega_0^{-1}\otimes\underline\omega_0\simeq\underline O_S$,
[formé des sections \inserted{locales} de $\mathscr J_0\otimes\underline\omega_0$ qui, pour
un \unsure{splittge|splitting} convenable de $\mathscr J_0$ \unsure{relmt} à $\underline\omega_0$, \inserted{(série de Hausdorff)} i.e.
\[
\mathscr J_0\simeq\underline\omega_0\oplus\underline O_S\oplus\underline\omega_0^{-1}
\]
\margin{splitting aussi donné par $\underline O_S\subset L_0$ correspond$^{\mathrm t}$ aux jets d'ordre $2$ de trivialisations \deleted{\unsure{canoniques}} ; ou ceux $\underline O_S\otimes\underline\omega$ \unsure{induis$^{\mathrm t}$} : $\omega_0\longrightarrow L_0\otimes\underline\omega_0$}
s'écrivent sous la forme
\[
\delta_2+0+1\quad\in\Gamma(\underline\omega_0^2+\underline\omega_0+\underline O_S)
\]

Les sections de ce sous-torseur correspondent
aux sections de $\mathbb P(\omega_0+\omega_0^{-1})$ qui sont disjointes
de la section canonique évidente de
$\omega_0\subset\omega_0+\omega_0^{-1}$, -- encore aux splittings
de $0\longrightarrow\omega_0\longrightarrow\omega_0+\omega_0^{-1}\longrightarrow\omega_0^{-1}\longrightarrow0$.
Il y en a une canonique (prenant pour celle-ci
celle de \unsure{Kampbell|Campbell}-Hausdorff) correspondant à la section $0+0+1$ de
$\mathscr J_0\otimes\underline\omega_0$. Elle est \uline{caractérisée} \inserted{parmi les conn. affines inv.} par la condition
que \deleted{\unsure{si $P$}}\inserted{loc. sur $S$ ($\beta$)} l'on puisse trouver sur
$P$ une \inserted{$G_a$-}réduction \inserted{invariante par translations} du groupe structural
à $G_a$, i.e. qu'on puisse trouver une
% margin note written sideways
\margin{$\beta$ Si on veut quelque chose de global sur $S$, il faut remplacer $G_a$ par $V(\underline\omega)$}
% archive sheet number: 138
\scanpage{140}
$\hookrightarrow$ extension sur $X$ invariante par translation
\[
0\mathbin{-}\unsure{\underline O_X|\underline t_X|\underline\omega_X}\mathbin{-}F\mathbin{-}\underline O_X\mathbin{-}0
\]
et un $G$-isom. $P\simeq\mathbb P(F)$ ($\simeq\overline{\mathbb V(\mathcal T)}$, où $\mathcal T$ est
le $G$-torseur sous $\underline O_X$, avec la $G$-structure
habituelle de $\underline O_X$, définie par l'extension \ldots)

[La \deleted{section}\inserted{horiz. $\sigma'$} de $P$ qui y correspond est
\unsure{partout|toute} disjointe de $\sigma$, donc correspond à
un splitting \inserted{(inv. par $H$)} $\mathscr J\simeq\underline\omega\oplus\underline O_X\oplus\underline\omega^{-1}$,
[qui d'ailleurs n'est autre que celui
déduit du splitting de Campbell-Hausdorff
sur l'alg. tronquée \ldots) avec $\underline O_X+\underline\omega^{-1}$ \inserted{et $\underline\omega^{-1}$} invar$^{\mathrm ts}$
pour la connexion ; et celle-ci induisant
la connexion \deleted{nulle} triviale sur le
quotient $\underline O_X$ \unsure{\ldots}] Pour les conn. affines
invariantes que l'on déduit par
la méthode de la ``conique'', il
y a $\hookrightarrow$ réduction \inserted{invariante} du groupe structural \deleted{au}
à un \inserted{un conjugué de,} \unsure{stabilis.} \unsure{sour|sous} de $G_0=G_t(\mathbb P(\underline\omega_0+\underline O_S))$.
Si on est sur un corps alg. clos (suffisant ici), \unsure{donc|alors}
la connexion proj. can. est caractérisée, parmi
celles qui sont invariantes, par les conditions
\deleted{de}\inserted{que} $P$ est \unsure{astreinte|admette} qu'une seule \deleted{\unsure{connexion|horizontale}}\inserted{section plate}
\deleted{\unsure{lin.}} invariante (\uline{NB} en car. $0$, \deleted{invariante}\inserted{plate}
? $\mid$ $\Rightarrow$ invariante), ttes les autres en
ont exactement deux \ldots.
% margin note written sideways
\margin{et si $T_0$ est le sous \unsure{tors.|tour} correspondant à $G_0$ (\unsure{uniqt déterminé|image déterminant}) on a $\unsure{\mathfrak P|\mathbb P}\simeq f^*(\overline{\mathcal T_0})=f^*(\overline{T_0})$ (\unsure{Algèbre projective|algébro-géomètre}), -- \unsure{conséquence} : c$^{\text{e}}$ section $1$, et la \deleted{\unsure{consid}}\unsure{restriction} \deleted{\unsure{elle}} à l'infini \deleted{\unsure{étant}} plate.}

\uline{NB} Si $H$ est propre sur $S$, fibres connexes (donc
$X/S$ un schéma en courbes elliptiques)
\uline{\unsure{tous|ttes}} les connexions projectives sur $X$
sont inv. par $H$. Ce qui caractérise donc
la conn. \inserted{proj.} canonique parmi \uline{t$^{\mathrm{tes}}$}
les connexions \ldots
% archive sheet number: 139
\scanpage{142}
1) \uline{Connexions sur torseurs affines}

$S/T$ schéma relatif \inserted{$(\unsure{g|q}:S\mathbin{-}T)$}, $E$ module qu. \unsure{ch.|coh.} sur $S$, d'où
\[
0\longrightarrow E\otimes\Omega^1_{S/T}\longrightarrow P^1_{S/T}(E)
\longrightarrow E\longrightarrow0 \tag*{(1)}
\]
extension de $\underline O_S$-modules. On a un splitting $g^{-1}(\underline O_T)$-linéaire
(mais pas $\underline O_S$-linéaire)
\[
d_E^1:E\longrightarrow P^1_{S/T}(E) \tag*{(2)}
\]
qui permet d'identifier $P^1_{S/T}(E)$ \uline{additiv$^{\mathrm{mt}}$} à $E\oplus E\otimes\Omega^1_{S/T}$, en
écrivant chaque \deleted{élément}\inserted{section} de $P^1_{S/T}(E)$
\inserted{(sur un ouvert)} de façon unique comme
\[
\zeta=\deleted{\unsure{\omega+}}\underset{\deleted{\unsure{E\otimes\Omega^1_{S/T}}}}{\,}
 d_E^{(1)}x+\omega\quad
 \left\{\begin{aligned}\omega&\in\Gamma(\ ,E\otimes\Omega^1_{S/T})\\
 x&\in\Gamma(\ ,E)\end{aligned}\right. \tag*{(3)}
\]
On a, si \deleted{\ensuremath{f}} $f$ est une section de $\underline O_S$
\[
f\zeta=f\omega+fd_E^{(1)}x=f\omega+
\underset{=\ d^1_{\unsure{\bullet}}(f)d^1_E(x)}{\underbrace{d_E^{(1)}(fx)}}+
\underset{-df}{[f-d_{\unsure{E|\bullet}}^{(1)}(f)]}d_E^1(x)
\]
i.e.
\[
\deleted{f(\omega+d_E^{(1)}x)=f\omega}\tag*{\deleted{(4)}}
\]
\[
f(d_E^1x+\omega)=d_E^1fx+(f\omega-x\otimes df).\tag*{(4)}
\]

Une \uline{connexion} \inserted{$\delta_E$} de $E$ rel. à $T$ est un splitting \inserted{$\underline O_S$-linéaire} de (1) ;
\deleted{\unsure{don|donc}} tant que splitting additif, compte tenu de
(2), il est donc donné par un hom additif
\[
d_E^0:E\longrightarrow E\otimes\Omega^1_{S/T} \tag*{(5)}
\]
par
\[
\delta_E x=d_E^1(x)-d_E^0(x), \tag*{(6)}
\]
il faut expliciter la condition sur $d_E^0$ de (5),
pour que $\delta_X$ défini par (6) soit $\underline O_S$-linéaire,
% archive sheet number: 141
\scanpage{143}
i.e. qu'on ait $f\delta x=\delta fx$, i.e. compte tenu de (4),
\[
d_E^1(fx)+(-fd_E^0x-x\otimes df)=d_E^1(fx)-d_E^0(fx),
\]
ce qui s'écrit
\[
d_E^0(fx)=fd_E^0(x)+x\otimes df \tag*{(7)}
\]
ou aussi, utilisant la structure de $\Omega^{\unsure{\cdot|\bullet}}_{S/T}$-module
à droite de $E\otimes\Omega^{\cdot}_{S/T}$ :
\[
d_E^0(\deleted{fx}xf)=d_E^0(x)f+x\,df.\tag*{(\deleted{\unsure{7}}7')}
\]

NB On dit souvent que $d_E^0$ (5) satisfaisant (7), est une connex. sur $E$ rel. à $T$.
\deleted{On sait alors que la \unsure{que}}

On sait alors qu'il $\exists$ prolongement unique de
$d_E^0$ en \deleted{une \unsure{dériv. extérieure|dérivation extérieure} \unsure{hom}} un endomorphisme
additif
\[
d_E^{\cdot}:E\otimes\Omega^{\cdot}_{S/T}\longrightarrow E\otimes\Omega^{\cdot}_{S/T} \tag*{(8)}
\]
qui \deleted{\unsure{est une anti}} satisfait à la relation suivante
rel. à la structure de $\Omega^{\cdot}_{S/T}$-module à droite
\[
d_E^{\cdot}(\varphi\omega)=d_E^{\cdot}(\varphi)\omega+
 (-1)^{\deg\varphi}\varphi\,d(\omega), \tag*{(9)}
\]
et ce $d_E^{\cdot}$ est un op de degré $1$. Pour que
l'on ait
\[
d_E^{\cdot}\circ d_E^{\cdot}=0, \tag*{(10)}
\]
il f. et s. que l'on ait
\[
\text{le composé }\ d_E^1d_E^0:E\longrightarrow E\otimes\Omega^1_{\unsure{S|E}/T}
\longrightarrow E\otimes\Omega^{\unsure{2}}_{\unsure{E|S}/T}\ \text{est nul.} \tag*{(10')}
\]

Le composé en question, considéré comme section
de $\operatorname{Hom}(E,E\otimes\Omega^2_{S/T})$, est appelé le \uline{tenseur courbure}
de la connexion. La connexion est dite \uline{intégrable}
si \unsure{le|la} \deleted{tenseur} courbure est \unsure{nul|nulle}, i.e. si \inserted{$d_E^{\cdot}$} (8) fait
% archive sheet number: 142
\scanpage{144}
de $E\otimes\Omega^{\unsure{\cdot|\bullet}}_{S/T}$ un complexe.

Supposons qu'il en soit ainsi. \deleted{\unsure{Rel. La $1^{\text{re}}$ tour d'un}}
\[
H^*_{\mathrm{DR}}(S,\underset{\substack{|\\\text{sous-entendu :}\\ E\text{ muni de }\delta_E\\
\text{i.e. de }d_E^0}}{E})=\mathbb H^*(S,E\otimes\Omega^{\cdot}_{S/T}). \tag*{(11)}
\]
On a une suite spectrale
\[
H^*_{\mathrm{DR}}(S,E)\Longleftarrow E_1^{pq}=
\deleted{\unsure{H^q(\mathrm{im}}}H^q(S,E\otimes\Omega^p_{S/T})) \tag*{(12)}
\]
associée à la filtration triviale du complexe $E\otimes\Omega^{\cdot}_{S/T}$.

On en \unsure{conclut|déduit} des \deleted{\unsure{aug}} homomorphismes
\[
H^*(\Gamma(S,E\otimes\Omega^{\cdot}_{S/T}))\Longrightarrow H^*_{\mathrm{DR}}(S,E) \tag*{(13)}
\]
qui sont des isom si p.ex. $S$ affine, et
\[
H^i_{\mathrm{DR}}(S,E)\longrightarrow H^i(S,E), \tag*{(14)}
\]
qui n'ont guère de raison d'être des isom. Ils sont
associés à l'hom de complexes évident (\unsure{troncature|troncation}) :
\[
E\otimes\Omega^{\cdot}_{S/T}\longrightarrow E. \tag*{(15)}
\]

\uline{NB} Pour $E$ module à conn. intégrable variable,
$H^*_{\mathrm{DR}}(S,E)$ est un $\unsure{\partial|\delta}$-foncteur en $E$ (On peut aussi
prendre la variante cat. dérivées). Les suites spectrales
et hom écrits plus haut sont fonctoriels en $E$.

\uline{Re NB} Images inverses \inserted{d'un Module à conn.} par

\diagram{\begin{tikzcd}[ampersand replacement=\&,row sep=small,column sep=small]
S\arrow[d]\& S'\arrow[l]\arrow[d]\\
T\& T'\arrow[l]
\end{tikzcd}}

carré commutatif :
tout suit \unsure{\ldots} !

En particulier ; \deleted{\unsure{img}} a) Image inverse par $T'\longrightarrow T$
(en prenant $S'=S\times_T T'$). b) Si $S\longrightarrow T'\longrightarrow T$, une
connexion rel. à $T$ en définit une rel. à $T'$, \deleted{et}
% archive sheet number: 143
\scanpage{145}
Nous allons étudier (14) pour $i=0,1$.

1) $H^0_{\mathrm{DR}}(S,E)$ est formée des sections \inserted{$x$} de $E$ qui sont
``horizontales pour la connexion'' $\delta$, ce qui \deleted{s'écrit par}\inserted{signifie par}
définition \deleted{$d_E^0x=0$}
\[
d_E^0(x)=0\qquad\text{i.e.}\qquad\delta_E x=d_E^1x \tag*{(16)}
\]
Alors (14) est l'injection canonique.

% slanted
\margin{Comparer Cat. \unsure{cof.|cat.|cf.} \unsure{additives|abéliennes|additions} 1.8.}

2) Soit $K^{\unsure{\cdot|\bullet}}:0\longrightarrow K^0\xrightarrow{d^0}K^1\xrightarrow{d^1}K^2\longrightarrow\cdots$
un complexe \deleted{\unsure{On}}\inserted{sur} \deleted{de faisceaux ab.}
\deleted{constate que $\mathbb H^1(S$}
de faisceaux abéliens sur un topos \inserted{$S$}. On constate que
$\mathbb H^1(S,K^{\cdot})$ s'interprète comme le groupe des
classes d'isomorphie \deleted{de} d'objets de la catégorie suivante \inserted{$C(K^{\cdot})$} :

$\operatorname{Ob} C(K^{\cdot})=$ triples $(P^0,\zeta)$, avec $P$ un torseur sous $K^0$,
$\zeta$ une section de $P^1=P\mathbin{\overset{K^0}{\times}}K^1$, telle que la section
\deleted{\unsure{canoni}} $d^1_*(\zeta)$ de $P^2=P\mathbin{\overset{K^0}{\times}}K^2\simeq
P^1\mathbin{\overset{K^1}{\times}}K^2$ soit la section
canonique, \deleted{\unsure{déi}} déduit du fait que $d^1d^0=0$.

$\mathrm{Fl}\ C(K^{\cdot})$ : on devine

\deleted{L'hom canonique $\mathbb H^1(S,K^{\cdot})\longrightarrow H^1(\unsure{S,}$}\inserted{Cette description de $\mathbb H^1(S,K^{\cdot})$ est compatible avec}
les hom. \unsure{en|entre} les complexes $K^{\cdot}$. Par suite,
l'hom \inserted{can.} \deleted{$\mathbb H^1(S,K^{\cdot}$}
\[
\mathbb H^1(S,K^{\cdot})\longrightarrow H^1(S,K^0)
\]
déduit de l'\deleted{\unsure{augmentation de $K$}}\inserted{hom troncature} $K^{\cdot}\longrightarrow K^0$ correspond au
foncteur $(P^0,\zeta)\longmapsto P^0$ ``oubli de $\zeta$''.

Noter que \deleted{la donnée, pour}\inserted{pour un} $P^0$ donné (torseur sous $K^0$) la
donnée d'une trivialisation $\zeta$ de $P^1$ équiv. à un isom
$P^1\simeq K^1$, donc les $\zeta$ correspondent \deleted{aux hom}
% archive sheet number: 144
\scanpage{146}
aux hom
\[
\bar d^0:P^0\longrightarrow K^1
\]
satisfaisant
\[
\bar d^0(xf)=\bar d^0(x)d^0(f)\qquad x\in P^0(\cdot),\ f\in K^0(\cdot).
\]
La condition mise ci-dessus sur $\zeta$ signifie simplement que $\bar d^0$ prend ses valeurs dans $Z^1K=\operatorname{Ker}(d^2)$.

Ceci s'applique notamment à la description
de $H^1_{\mathrm{DR}}(S,E)$, comme classifiant les objets suivants.

Un torseur $P$ sous $E$, un hom de faisceaux
d'ensembles
\[
d_P^{\deleted{\unsure{0}}}:P\longrightarrow\deleted{\unsure{Z^1(|Z(}}E\otimes\Omega^1_{S/T}\deleted{\unsure{)}}, \tag*{(17)}
\]
satisfaisant les conditions
\[
\left\{\begin{aligned}
d_P^{\deleted{\unsure{0}}}(xf)&=d_P^{\deleted{\unsure{0}}}(x)d_E^0(f)
 &x&\in P(\cdot),\ f\in\underline O_S(\cdot).\\
d_P(P)&\subset Z^1(E\otimes\Omega^{\unsure{1|\cdot|{}}}_{S/T})\ \text{i.e.}\ d_E^1d_P=0
\end{aligned}\right. \tag*{(18)}
\]
Un tel $d_P$ sur un $P$ s'appellera une \uline{connexion} \inserted{\uline{intégrable}} \uline{sur
le torseur} $P$, \uline{compatible avec la connexion} \inserted{\uline{intégrable}} $d_E^0$ \uline{sur} $E$.

\uline{NB} Si $d_E^0$ est une connexion sur $E$ (pas néc. intégrable),
\deleted{on peut \unsure{se pro}} on peut se proposer de trouver, pour un torseur $P$ sous $E$,
quelles sont les connexions compatibles avec $d_E^0$, au sens
général \deleted{\unsure{du ca}} de données de recollement infinitésimal
d'ordre $1$. Il est clair que si $S$ est affine, il existe
de telles connexions, donc il en existe tjrs localemt,
et que les connexions forment un torseur sous
$E\otimes\Omega^1$, qui en fait est isomorphe à $P\mathbin{\overset E{\times}}(E\otimes\Omega^1)$
(grâce à $d_E^0:E\longrightarrow E\otimes\Omega^1$). Donc les connexions
% archive sheet number: 145
\scanpage{147}
On trouve en particulier une interprétation de
$H^1_{\mathrm{DR}}(S,E)$, comme classifiant les torseurs $P$ sous $E$,
munis d'un
% archive sheet number: 146
\scanpage{148}
s'interprètent comme les hom (17) satisfaisant
la \uline{première relation} (18). \deleted{\unsure{C'est lorsqu|C'est à lorsqu}} C'est la
deuxième relation qui précise que l'on se restreint
aux connexions intégrables sur $P$.

Les \deleted{faisceau des} automorphismes de $P$ \uline{compatibles avec} $d_P$
\deleted{2) \uline{Stratifications sur torseurs affines}}
est \deleted{\unsure{form}} sous-gpe du gpe $\Gamma(S,E)$ des automorphismes
de $P$ formé des \deleted{$x\in\Gamma(S,E$} sections \uline{horizontales} de $E$,
i.e. $H_{\mathrm{DR}}^0(S,E)$.

Les descriptions qui précèdent sont \inserted{fonctorielles en $E$, et} compatibles
aux images inverses envisagés (sic) plus haut.

2) \uline{Stratifications sur torseurs affines}

Supposons donnée sur $E$ une stratification rel. à $T$.
Alors $E$ définit un faisceau \inserted{de modules $E_{\mathrm{strat}}$} du site \unsure{stratifiant|stratifié}
$(S/T)_{\mathrm{strat}}$, donc on peut former
\[
H^*_{\mathrm{strat}}(S,E)\overset{\unsure{\mathrm{dfn}|\mathrm{déf}}}{=}H^*((S/T)_{\mathrm{strat}},E_{\mathrm{strat}}). \tag*{(19)}
\]
\deleted{La stratification de $E$}
\deleted{On a défini dans [\quad] un hom naturel}

Une donnée de recollement \unsure{d'infinitésimale|infinitésimale} d'ordre $\infty$ sur
$E$ revient à la donnée d'un homomorphismes
\margin{On dit parfois que $\delta_E^\infty$ est la connexion.}
\[
\delta_E^\infty:E\longrightarrow P^\infty_{S/T}(E) \tag*{(20)}
\]
satisfaisant la condition
\[
\varepsilon_E\delta_E^\infty:E\longrightarrow P^\infty_{S/T}(E)
\xrightarrow{\varepsilon_E} E\quad\text{est l'identité.} \tag*{(21)}
\]
\margin{$P^\infty_{S/T}(E)^+$, idéal d'augmentation}
\inserted{On préfèrera $\deleted{\unsure{d_E}}d_E^+:E\longrightarrow P^\infty_{S/T}(E)^+,\quad
d_E^+(x)=d^\infty_{\deleted{\unsure{E}}}(x)-\delta_E^\infty(x)$}
\deleted{(21) La condition pour la donnée de recollement inf. soit
une stratification est une condition sur $\delta_E$ qui}
% slanted
\margin{plus bas, ds n° 3.}
% archive sheet number: 147
\scanpage{149}
\deleted{\unsure{indique que|implique que}} Le composé de (\unsure{22|21}) avec \deleted{$P^\infty_{S/T}$}
\[
P^\infty_{S/T}(E)^+\longrightarrow P^1_{S/T}(E)^+\simeq E\otimes\Omega^1_{S/T},\ \deleted{\text{\unsure{est}}}
\]
noté $d_E^0$
\[
d_E^0:E\longrightarrow E\otimes\Omega^1_{S/T},\quad d_E^0=d_E^+\bmod\ldots \tag*{(23)}
\]
est l'hom (5) associé à la connexion de $E$ définie
par la donnée de recollement. \deleted{Si} \uline{On vérifie que si
la donnée de recollement} \deleted{définie par $\delta_E$} \uline{est une stratification},
alors la connexion associée est intégrable, i.e. (23)
donne lieu à un complexe de De Rham.

On sait \unsure{ds|alors|bien} [\quad] définir des hom canoniques
\[
H^*_{\mathrm{strat}}(S,E)\longrightarrow H^*_{\mathrm{DR}}(S,E), \tag*{(24)}
\]
$\unsure{\partial|\delta}$-fonctoriels en $F$, compatibles au changemt de base.

Nous voulons une interprétation de $H^1_{\mathrm{strat}}(S,E)$
analogue à celle donnée plus haut pour $H^1_{\mathrm{DR}}(S,E)$.

\deleted{Pour ceci, \deleted{\unsure{posons}} Soit, \inserted{pour tt $\unsure{i\in\mathbb N|i\geq0}$,} $C^i(E)$
\deleted{l'image inverse de $E$} \deleted{\unsure{sur le}} le faisceau sur
$S_{\mathrm{Zar}}$ \deleted{image inverse du} complété formel le long de la diagonale de
$(S/T)^{i+1}$ de l'image inverse de $E$ par (disons) $\operatorname{pr}_{\unsure{1|i}}$.
\deleted{La stratification de}}

\[
\deleted{C^0(E)=E,\quad C^1(E)=P^\infty_{S/T}(E).}\tag*{\deleted{(24)}}
\]

\deleted{La stratification de $E$ met sur $C^{\unsure{\cdot|\bullet}}(E)$ une structure
de faisceau cosimplicial [en fait, un co-module
cosimplicial sur $C^{\cdot}(\underline O_S)$] :}

\diagram{\begin{tikzcd}[ampersand replacement=\&,column sep=small,row sep=small,execute at end picture={\draw[thin] (current bounding box.south west) -- (current bounding box.north east);\draw[thin] (current bounding box.north west) -- (current bounding box.south east);}]
\deleted{(25)}\&\deleted{\deleted{\unsure{E^0|E}}}\arrow[r,shift left,"\deleted{\delta_E^\infty}"]\arrow[r,shift right,swap,"\deleted{d_E^\infty}"]\&
\deleted{\deleted{P_{S/T}^\infty(E)}}\&\deleted{C^0(E)}\arrow[r,shift left]\arrow[r,shift right]\&
\deleted{C^1(E)}\arrow[r,shift left=1.2ex]\arrow[r]\arrow[r,shift right=1.2ex]\&\deleted{C^2(E)\ldots}\\
\&\&\&\deleted{E}\arrow[u,equal]\arrow[r,shift left,"\deleted{\delta_E^\infty}"]\arrow[r,shift right,swap,"\deleted{d_E^\infty}"]\&
\deleted{P^\infty_{S/T}(E)}\&\deleted{,}
\end{tikzcd}}
% archive sheet number: 148
\scanpage{150}
Or un torseur sur $S_{\mathrm{strat}}$ sous $E_{\mathrm{strat}}$ peut être interprété
comme la donnée, pour tout objet \deleted{de} \inserted{$(U,U')$ de} $S_{\mathrm{strat}}$, d'un
torseur sous $(E_{\mathrm{strat}})_{(U,U')}=E_{U'}$, \deleted{\unsure{est}} qui varie de façon
contravariante en $(U,U')$. On voit que ceci équivaut
à la donnée d'un torseur ordinaire $P$ sous $E$, et
d'une \uline{stratification} sur $P$, compatible avec celle de $E$.

\diagram{\begin{tikzcd}[ampersand replacement=\&,row sep=tiny,column sep=small]
U\arrow[r,phantom,"\subset" description]\arrow[d,phantom,"{\rotatebox{-90}{$\subset$}}" description]\& U'\arrow[ddl]\\
S\arrow[d]\&\\ T\&
\end{tikzcd}}

\deleted{Nous regardons} Nous regardons $P^\infty_{S/T}(E)$ comme un \uline{préfaisceau} \inserted{A.R.}
sur $S$, et comme tel il est muni d'une stratification
canonique (\unsure{indépend$^{\mathrm{te}}$|indépendante} de celle de $E$, cf [\quad]).
Nous \unsure{obtenons|avons} la \deleted{\unsure{hom ca}} suite exacte
\[
0\longrightarrow E\xrightarrow{\delta=\delta_E^\infty}P^\infty_{S/T}(E)
\xrightarrow{p}P^\infty_{S/T}(E)/\delta(E)\longrightarrow0, \tag*{(25)}
\]
où le dernier terme est défini comme le \unsure{syst.|sup|support} projectif des \deleted{\unsure{coimages}}\inserted{conoyaux} des composés
$E\xrightarrow{\delta_E^\infty}P^\infty_{S/T}(E)\longrightarrow P^n_{S/T}(E)$.

\deleted{\unsure{Comme les hom}} Notons que l'\uline{hom. $\delta$ est compatible avec les
stratifications}. La donnée d'un torseur stratifié $P_{\mathrm{strat}}$
sous $E$ s'identifie alors à la donnée d'un
torseur stratifié sous le préfaisceau $P^\infty_{S/T}(E)$, et
d'une trivialisation \deleted{\unsure{compat}} du \deleted{\unsure{pr}} torseur stratifié associé
de groupe $P^\infty_{S/T}(E)$. Or on constate qu'un torseur
stratifié sous $P^\infty_{S/T}(E)$, c'est ``la $=$ chose'' qu'un
torseur ordinaire \inserted{$P$} sous $E$, (cf. [\quad] 6.4), à un
tel $P$ étant associé le préfaisceau \inserted{$P^\infty_{S/T}(P)$} des \deleted{\unsure{parties}}
jets de $P$ sur $S/T$ d'ordre \unsure{$-$|$n$|$\infty$} (dont \unsure{$-$|$\infty$}). Alors
la trivialisation $P^\infty_{S/T}(P)\longrightarrow P^\infty_{S/T}(E)/\delta E$ revient à la
% archive sheet number: 149
